\section{Alcance de la abstracción poblacional y estabilidad de la red de precios}
\label{sec:alcance}

\subsection{Dónde vale tratar la flota como población}

El marco de juegos de población que sostiene \cref{obs:smith} y
\cref{th:smallgain} se enuncia para poblaciones cuyo número de agentes es
«fijo y típicamente grande», y en el que «las identidades individuales de los
agentes no son relevantes, solo sus estrategias»
\cite{martinezPiazuelo2025gne}. Ambas condiciones merecen escrutinio aquí,
porque este trabajo trata explícitamente con AMR heterogéneos y flotas de una
decena.

\begin{observacion}[Tres regímenes de validez]
\label{obs:alcance-poblacional}
La abstracción poblacional no se aplica de manera uniforme a las tres capas:
\begin{enumerate}[label=(\alph*),leftmargin=2.2em,topsep=0.2em,itemsep=0.15em]
  \item \textbf{Reparto de esfuerzo (exacta).} La masa que se distribuye es el
  esfuerzo de contacto $\bm\lambda$, magnitud continua sobre un conjunto
  convexo. No hay agentes que contar: la abstracción es exacta, no aproximada.
  \item \textbf{Precios y congestión (exacta bajo convexidad).} El objeto es
  una masa continua sobre recursos compartidos, y el análisis es el de una
  desigualdad variacional.
  \item \textbf{Reclutamiento (aproximada).} Aquí sí hay agentes discretos,
  heterogéneos y pocos. La identidad \emph{sí} importa: la capacidad efectiva
  $e_{ik}$ distingue a un AMR de otro, y un miembro de una coalición de cuatro
  aporta una fracción no despreciable del \emph{wrench}. La hipótesis de
  influencia individual despreciable falla por construcción.
\end{enumerate}
\end{observacion}

Que falle en (c) no invalida el trabajo: obliga a decir qué se conserva. Se
conserva todo lo demostrado sobre el juego \emph{atómico} —
\cref{th:potencial}, \cref{co:terminacion}, \cref{th:noequiv},
\cref{pr:jerarquia}, \cref{th:descomposicion}—, porque ninguno de esos
resultados invoca población grande: están enunciados sobre perfiles discretos y
desviaciones factibles. Lo que queda condicionado es el uso de las dinámicas
poblacionales como \emph{descripción del transitorio} del reclutamiento.

\begin{observacion}[El ruido de exploración no es un parámetro libre]
\label{obs:ruido}
En el límite de población grande, el proceso de revisión estocástico converge a
la dinámica determinista con fluctuaciones que decrecen con la masa. A $N$
finito esas fluctuaciones persisten, y su papel es doble: son el error respecto
al límite determinista y son, a la vez, el mecanismo que permite escapar de
terminales de mejor respuesta. De ahí una regla de diseño coherente con
\cref{sec:nhop}: las capas que \emph{rastrean} dentro de una cuenca —reparto de
esfuerzo, precios— deben operar con masa efectiva alta y comportamiento
determinista; la capa que \emph{selecciona} entre cuencas —reclutamiento—
requiere fluctuación, y por eso el logit se emplea como perturbación de
exploración y no como descripción del reposo. La temperatura no es un parámetro
libre: compite con $h_c^\star$ como mecanismo de escape.
\end{observacion}

\subsection{Por qué el reparto es exacto: el enfoque por analogía}

El apartado (a) afirmaba exactitud sin justificarla más allá de la continuidad.
La razón es más precisa y tiene nombre: la literatura de juegos de población
distingue el \emph{enfoque natural} —cuando el sistema consta genuinamente de
muchos agentes que eligen de forma esporádica y con racionalidad acotada— del
\emph{enfoque por analogía}, en el que las variables de control de un sistema con
restricción de tipo simplex se interpretan como el estado social de una
población, aunque no haya población literal \cite{martinezPiazuelo2025gne}.

\begin{observacion}[El reparto en cuotas es un simplex, y eso lo da todo]
\label{obs:analogia}
Escríbase el reparto en \emph{cuotas} $\alpha_i\ge0$ con
$\sum_{i\in\Ccal_\ell}\alpha_i=1$, en lugar de en esfuerzos absolutos. Ese
conjunto es exactamente un simplex, que es la estructura que el enfoque por
analogía requiere. Y entonces las dos invariancias que las dinámicas evolutivas
preservan por construcción —no negatividad y conservación de masa— \emph{son} las
dos restricciones físicas del problema: la no negatividad es la unilateralidad
de \cref{pr:notiron}, es decir que ningún AMR tira de la carga; y la
conservación de masa es que la demanda queda íntegramente cubierta.

La comprobación distingue los dos métodos donde importa. Con cuatro AMR y un
coste cuyo mínimo sobre el plano $\sum\alpha_i=1$ tiene una componente negativa
—$(-2;1;1;1)$—, el gradiente simple converge a ese punto: asigna cuota $-2$, que
físicamente es un AMR \emph{tirando} de la carga. La dinámica evolutiva se
detiene exactamente en $\alpha_0=0$ y conserva la masa a $1{,}000000$.

La consecuencia práctica es que no hace falta paso de proyección, ni comprobar a
posteriori que los esfuerzos son no negativos: la estructura de la dinámica lo
garantiza. Lo que en otra formulación sería una restricción a imponer, aquí es
una propiedad del método.
\end{observacion}

\subsection{Dos formas de imponer una restricción}

Queda una distinción de diseño que este trabajo emplea en dos sitios sin haberla
nombrado. Ante una restricción acoplada caben dos estrategias
\cite{martinezPiazuelo2025gne}: la \emph{prescriptiva}, que rediseña el protocolo
de revisión para que la restricción no pueda violarse —multiplicando la tasa de
cambio hacia una opción por el margen que le queda—, y la \emph{no
prescriptiva}, que deja el protocolo intacto y añade un multiplicador, es decir,
un precio.

\begin{observacion}[El precio no protege durante el transitorio]
\label{obs:prescriptivo}
La diferencia no es de matiz. Sobre una población con tres estrategias, donde la
más atractiva tiene capacidad $0{,}35$, el mecanismo de precios alcanza una
ocupación máxima de $0{,}753$ —un $115\%$ por encima de la capacidad— y permanece
por encima de ella durante el $51\%$ del horizonte, antes de converger a
$0{,}3495$. El protocolo prescriptivo alcanza $0{,}3500$ y \emph{nunca} la
excede.

Y no es cuestión de sintonía: subir la ganancia del ajuste de precios reduce el
exceso pero no lo anula —$+0{,}564$ con ganancia $2$, $+0{,}403$ con $8$,
$+0{,}227$ con $40$ y $+0{,}114$ con $200$—. Un precio impone la restricción en
el equilibrio; un protocolo la impone en todo instante.
\end{observacion}

\begin{observacion}[El criterio que ordena la arquitectura]
\label{obs:criterio-restricciones}
De ahí se sigue el principio que este trabajo ya aplica sin enunciarlo: \emph{una
restricción cuya violación es catastrófica debe imponerse de forma prescriptiva;
una cuya violación es meramente costosa puede imponerse con un precio}. Por eso
la separación entre vecinos se trata con barreras —\cref{pr:hocbf} y
\cref{th:astop-cuerpo}, que dan invariancia en todo instante— y la congestión de
corredores con peajes (\cref{th:multiclase}), que solo la internalizan en el
equilibrio.

Conviene por tanto no leer el peaje multiclase como un mecanismo de capacidad.
No lo es: cobra la demora que una coalición ancha impone a las demás, no impide
que el corredor se sobrecargue. Quien impide la sobrecarga es la barrera, y el
peaje solo desincentiva llegar a ella. Confundir ambos papeles sería atribuir al
precio una garantía que \cref{obs:prescriptivo} muestra que no puede dar.
\end{observacion}

\subsection{Qué preserva la composición y qué no}

\Cref{th:smallgain} compone lazos disipativos sin decir qué operaciones conservan
esa propiedad. Conviene enunciar las dos que este trabajo utiliza
\cite{martinezPiazuelo2025gne}.

\begin{proposicion}[Dos composiciones que conservan la disipatividad]
\label{pr:composiciones}
Sea $\Sigma$ un sistema $\delta$-disipativo caracterizado por $(\sigma,\Gamma)$.
\begin{enumerate}[label=(\roman*),leftmargin=2.2em,topsep=0.2em,itemsep=0.15em]
  \item \emph{Ganancia en cascada}: multiplicar la salida por $\kappa\neq0$
  conserva la $\delta$-disipatividad, con la matriz de suministro reescalada por
  $\kappa$ en los bloques cruzados y por $\kappa^{2}$ en el segundo diagonal. En
  particular, con $\kappa=-1$ un sistema $\delta$-pasivo se vuelve
  $\delta$-antipasivo con el mismo $\sigma$, y recíprocamente.
  \item \emph{Encapsulado sin memoria}: envolver un sistema escalar
  $\delta$-pasivo con una función suave $\varphi$ conserva la $\delta$-pasividad
  si $\varphi$ es afín, o bien si el sistema es positivo, arranca en el ortante
  no negativo y $\varphi$ es \emph{convexa}. La función de $\delta$-almacenamiento
  del encapsulado es $\bar S(z,\bar u)=S(z,\varphi(\bar u))$.
\end{enumerate}
\end{proposicion}

\begin{observacion}[Por qué la realimentación negativa es lo que hace funcionar el lazo]
\label{obs:signo-menos}
El apartado (i) explica algo que la arquitectura usa sin justificar. La condición
de estabilidad exige que la dinámica de decisiones sea $\delta$-pasiva y la de
pagos $\delta$-antipasiva: son propiedades opuestas y parecería que hay que
diseñarlas por separado. No hace falta. El signo negativo de la realimentación
—que el precio \emph{penalice} en lugar de premiar la congestión— convierte una
en la otra conservando el mismo $\sigma$. Que el ajuste \eqref{eq:tatonnement}
suba el precio cuando hay exceso de demanda no es una convención económica: es lo
que hace que el par cumpla la condición de interconexión.
\end{observacion}

\begin{observacion}[La convexidad es condición suficiente, no necesaria]
\label{obs:convexidad-encapsulado}
El apartado (ii) da una regla de diseño: toda no linealidad insertada en el lazo
de precios —una función de urgencia, una ponderación de prioridad, una
saturación— debe ser convexa o afín para que el certificado de composición se
conserve. Los precios de este trabajo son no negativos y las funciones de coste
convexas, de modo que la hipótesis se cumple.

Conviene no exagerar el alcance, y la comprobación numérica obliga a ello.
Ensayando el lazo encapsulado con envolturas convexas, afines y cóncavas, y con
dos ganancias distintas, las variantes cóncavas resultaron en su mayoría
igualmente estables. Lo que se pierde al salir de la clase convexa no es la
estabilidad sino el \emph{certificado}: se deja de poder afirmarla por
composición y habría que demostrarla caso por caso. La única patología observada
fue con una envoltura de derivada no acotada en el origen, que produjo
castañeteo en torno al reposo sin llegar a divergir.
\end{observacion}

\subsection{Restringir a quién se puede cambiar}

En este trabajo un AMR no puede reasignarse a cualquier carga: solo a las que
alcanza dentro de su bola de $n$ saltos. Eso es una \emph{interacción restringida
por estrategia}, que se modela con un grafo cuyos nodos son las estrategias y
cuyas aristas indican qué intercambios están permitidos, y que se emplea
precisamente para representar estructuras de información local
\cite{quijano2017role}. Conviene saber qué cuesta esa restricción.

\begin{observacion}[Basta con que el grafo de estrategias sea conexo]
\label{obs:grafo-estrategias}
Sobre un juego de cuatro cargas con equilibrio de Nash
$(0{,}4951;\ 0{,}3689;\ 0{,}1359;\ 0)$, arrancando desde un reparto muy
desequilibrado:
\begin{itemize}[leftmargin=1.9em,topsep=0.2em,itemsep=0.15em]
  \item grafo completo: converge exactamente al equilibrio;
  \item grafo camino —conexo pero con solo tres aristas de seis—: converge
  \emph{exactamente al mismo} equilibrio, con error nulo;
  \item grafo con dos componentes: converge a
  $(0{,}4667;\ 0{,}3333;\ 0{,}2;\ 0)$, a distancia $0{,}079$ del equilibrio, y la
  masa queda atrapada en cada componente con el reparto inicial, $0{,}80$ y
  $0{,}20$.
\end{itemize}

La lectura es tranquilizadora y precisa a la vez: \emph{la escasez de aristas no
desplaza el equilibrio; la desconexión sí}. Mientras la reachability entre cargas
se mantenga conexa, restringir los cambios a $n$ saltos no cuesta exactitud —solo
velocidad, por el diámetro—. Si se fragmenta, cada componente conserva su masa
por separado y el reposo pasa a depender de cómo estuviera repartida la flota al
fragmentarse, que es una dependencia de condición inicial que ningún resultado de
este trabajo contempla.

Esto refuerza la condición de carrera de \cref{obs:carrera} desde otro ángulo: no
solo importa que el reclutamiento converja antes de la separación, sino que
mientras esté activo el grafo de estrategias no se parta.
\end{observacion}

\subsection{Cuántos saltos hacen falta, y qué compra la información local}

La pregunta «¿cuántos saltos necesita el mecanismo?» no tiene respuesta
universal, pero sí la tiene para el subproblema definido positivo que resuelve
cada rama cuadrática del reparto, y la respuesta pasa por el condicionamiento.

\begin{teorema}[Inversión local y alcance de dependencia]
\label{th:saltos}
Sea $K=K^{\!\top}\succ0$ \cite{mayorga2026m1} con $\mu I\preceq K\preceq LI$, cuyo grafo de
dispersión conecta dos variables cuando su bloque es no nulo, de modo que un
producto $Kz$ cuesta una ronda de intercambio. Con $\gamma=2/(\mu+L)$,
$\rho=(L-\mu)/(L+\mu)<1$ y $p_h(K)=\gamma\sum_{j=0}^{h}(I-\gamma K)^{j}$, que
usa exactamente $h$ saltos de información,
\[
  \lVert K^{-1}-p_h(K)\rVert_2\ \le\ \frac{\rho^{\,h+1}}{\mu}.
\]
\end{teorema}

\begin{observacion}[La cota es exacta, y el condicionamiento manda]
\label{obs:saltos-num}
Sobre matrices de anillo con autovalores en $[1,\kappa]$, el error medido
coincide con la cota en todos los casos hasta precisión de máquina: con
$\kappa=2$, $3{,}70\times10^{-2}$ en $h=2$ y $5{,}65\times10^{-6}$ en $h=10$;
con $\kappa=50$, $0{,}887$ en $h=2$ y todavía $0{,}194$ en $h=40$.

Los saltos necesarios para un error de $10^{-3}$ son $6$, $34$ y $172$ para
$\kappa=2$, $10$ y $50$. Esto enlaza dos resultados que estaban separados: la
degradación de autoridad de \cref{th:degradacion} empeora el condicionamiento
del reparto, y por este teorema \emph{un reparto mal condicionado exige más
saltos para la misma precisión}. La autoridad física y la profundidad de
comunicación no son recursos independientes; el segundo compensa al primero, y
a un precio geométrico en $\rho$.
\end{observacion}

Lo anterior dice cuánta información hace falta. Lo siguiente dice qué se puede
certificar con la que hay.

\begin{teorema}[Intervalo de optimalidad desde información aproximada]
\label{th:intervalo}
Sea una rama con conjunto cerrado y convexo, coste $\mu$-fuertemente convexo
$\Phi$ y candidato factible $z$. Si se dispone de una aproximación local
$\hat g$ del gradiente con $\lVert\hat g-\nabla\Phi(z)\rVert\le\varepsilon_g$ y
un residual de estacionariedad $\varepsilon_r$ \cite{mayorga2026m1}, entonces el mínimo $\Phi^\star$
satisface
\[
  \Phi(z)-\frac{(\varepsilon_r+\varepsilon_g)^{2}}{2\mu}\ \le\ \Phi^\star\ \le\ \Phi(z).
\]
\end{teorema}

\begin{observacion}[Un certificado que se calcula con lo que se tiene]
\label{obs:intervalo-num}
Con $\mu=2$ y errores de gradiente de $0$, $0{,}1$ y $0{,}5$, el intervalo
contiene al óptimo en los tres casos, y su anchura crece con el error:
$0{,}249$, $0{,}062$ y $0{,}660$ respectivamente para los candidatos ensayados.

Su valor para este trabajo es que convierte en cantidad una cosa que
\cref{th:descomposicion} solo nombraba: la pérdida atribuible a información
imperfecta. Cada AMR puede acotar, con su gradiente local aproximado, cuánto
pierde su decisión respecto de la óptima; y la regla de aceptación pasa a ser
verificable —aceptar una decisión solo cuando su intervalo la distingue de las
alternativas—, que es la forma correcta de leer la selección con certificados
locales que sostiene \cref{th:testigo}.
\end{observacion}

\subsection{Cuánto se aleja este juego de tener equilibrio único}

Hay una magnitud que resume por qué el reclutamiento es la capa difícil y las
otras no.

\begin{definicion}[Déficit de contractividad]
\label{def:deficit}
Un juego de población $\Gcal=(\Pcal,\Xcal,f)$ es \emph{contractivo} si
$(x-\tilde x)^{\!\top}(f(x)-f(\tilde x))\le0$ para todo par de distribuciones. Un
juego cualquiera se caracteriza por su \emph{déficit}
\[
  \nu^\star=\inf\Bigl\{\nu:\ (x-\tilde x)^{\!\top}\bigl(f(x)-f(\tilde x)\bigr)
  \le\nu\,\lVert x-\tilde x\rVert^{2}\Bigr\},
\]
que es no positivo exactamente cuando el juego es contractivo
\cite{martinezPiazuelo2025gne}.
\end{definicion}

\begin{observacion}[Medido sobre las tres capas]
\label{obs:deficit-capas}
El déficit separa las capas de este trabajo con nitidez. Muestreando sobre el
simplex:
\begin{itemize}[leftmargin=1.9em,topsep=0.2em,itemsep=0.15em]
  \item reparto de esfuerzo con coste convexo: $\nu^\star=-2{,}42$, contractivo;
  \item congestión con coste creciente en la propia masa: $\nu^\star=-0{,}90$,
  contractivo;
  \item formación de coalición con sinergia $c$ —el pago de unirse \emph{crece}
  con cuántos se han unido—: $\nu^\star=-0{,}13$ para $c=0$, y $+0{,}20$,
  $+0{,}87$ y $+1{,}87$ para $c=0{,}5$, $1{,}5$ y $3$.
\end{itemize}

La lectura es que \emph{la sinergia es exactamente lo que destruye la
contractividad}, y el déficit crece de forma aproximadamente lineal con ella. Las
dos primeras capas son contractivas porque sus externalidades son negativas —más
esfuerzo cuesta más, más tráfico retrasa más—; la tercera no lo es porque la suya
es positiva: que otro se una \emph{mejora} el pago de estar dentro.
\end{observacion}

\begin{observacion}[Es la caza del ciervo, medida]
\label{obs:deficit-ciervo}
Esto da una segunda lectura de \cref{pr:caza-ciervo}. Un juego contractivo tiene
equilibrio único; la multiplicidad que allí se demuestra —«todos se unen» y
«nadie se une» conviviendo— solo es posible con déficit positivo. Coordinación
frágil y déficit de contractividad no son dos problemas: son el mismo, medido
desde la teoría de juegos y desde el análisis de dinámicas.

Y sitúa el papel del margen de \cref{th:smallgain}: lo que ese margen debe
absorber es precisamente $\nu^\star$. Un sistema cuya sinergia de coalición sea
alta necesita más margen disipativo en el resto de las capas, o dicho al revés,
la fragilidad de la coordinación se paga en la sintonía de todo lo demás.
\end{observacion}

\subsection{Qué exige de verdad difundir el precio}

Este trabajo afirma que los precios se difunden «por el mismo consenso» y remite a
\cref{th:conmutada}, cuya hipótesis es que la unión de topologías contenga un
árbol generador. La construcción distribuida de multiplicadores exige más que
eso, y conviene enunciarlo porque el modo de fallo es engañoso.

\begin{supuesto}[Grafo equilibrado en pesos]
\label{s:equilibrado}
El digrafo de comunicación entre proveedores de pago es fuertemente conexo y
\emph{equilibrado en pesos}: en cada nodo, la suma de pesos entrantes iguala a la
de salientes \cite{martinezPiazuelo2025gne}.
\end{supuesto}

\begin{observacion}[Sin equilibrio de pesos hay acuerdo, pero es el acuerdo
equivocado]
\label{obs:equilibrado}
La construcción asigna a cada población un proveedor de pago que solo observa su
distribución local y su parte de la restricción, y difunde con una estructura
proporcional-integral de \emph{dos} estados por nodo. El segundo estado integra
el desacuerdo, y es lo que lleva a todos al mismo multiplicador sin que ninguno
conozca la restricción completa.

Lo que ocurre al violar \cref{s:equilibrado} no es divergencia. Sobre un lazo
cerrado de cuatro nodos con restricción acoplada $\sum_k A_kx_k=b$ y $b=3$:
\begin{itemize}[leftmargin=1.9em,topsep=0.2em,itemsep=0.15em]
  \item con anillo no dirigido, equilibrado, las decisiones convergen
  \emph{exactamente} al óptimo centralizado, el multiplicador a $1{,}2444=\lambda^\star$
  y el error de restricción es $0{,}00000$;
  \item con un digrafo fuertemente conexo pero no equilibrado, los nodos
  \emph{también} alcanzan acuerdo —dispersión $9\times10^{-14}$— pero sobre el
  valor equivocado, $1{,}4$ en lugar de $1{,}2444$, y la restricción queda
  violada en $-0{,}583$: se entrega $2{,}417$ donde se pedía $3$.
\end{itemize}

Ese es el peligro: el sistema no avisa. Todos los nodos coinciden, el ajuste
reposa y nada indica avería, pero el multiplicador acordado no es el que vacía el
mercado. Un diagnóstico basado en la dispersión entre nodos —que es lo natural—
no detectaría nada. La comprobación pertinente no es que los precios coincidan,
sino que el residuo agregado de la restricción se anule.
\end{observacion}

\begin{observacion}[Dos limitaciones más de la construcción]
\label{obs:multiplicadores-limites}
Conviene declarar otras dos. Primera, la construcción distribuida está
desarrollada para restricciones de \emph{igualdad}; las de desigualdad —como la
capacidad de un corredor— requieren proyección adicional sobre el ortante no
negativo, y la garantía de composición debe rehacerse en ese caso. Segunda, el
coste de estado es el doble del que sugiere la intuición: dos variables por nodo
y por restricción, no una. En una flota con varias cargas activas y varios
corredores compartidos, ese factor dos multiplica el estado que cada AMR debe
mantener y difundir, y entra en la contabilidad de \cref{def:complejidad}.
\end{observacion}

\figLazo

\subsection{Estabilidad de la red de precios}

El trabajo introduce precios en tres lugares —\emph{wrench}, congestión y
seguridad (\cref{th:precio-seguridad})— y hasta aquí ha tratado cada lazo por
separado. Con varias cargas activas simultáneamente, los lazos se acoplan: subir
el precio de una carga atrae AMR que otra pierde. Esa interacción puede
desestabilizar el ajuste, y conviene enunciar cuándo no lo hace.

Sea $p=(p_\ell)_{\ell\in\Kcal}$ y el ajuste tipo \emph{tâtonnement}
\begin{equation}
  \dot p_\ell=\gamma\bigl(D_\ell-S_\ell(p)\bigr),
  \label{eq:tatonnement}
\end{equation}
con $S_\ell$ la capacidad efectiva que la carga $\ell$ atrae al precio vigente.
Sea $E$ la matriz de elasticidades cruzadas, $E_{\ell\ell'}=\partial S_\ell/\partial p_{\ell'}$.

\begin{proposicion}[Dominancia diagonal y estabilidad de precios]
\label{pr:dominancia}
Si en el punto de equilibrio la matriz $E$ satisface
\[
  E_{\ell\ell}>0
  \quad\text{y}\quad
  E_{\ell\ell}>\sum_{\ell'\neq\ell}\bigl|E_{\ell\ell'}\bigr|
  \qquad\text{para todo }\ell,
\]
entonces la linealización de \eqref{eq:tatonnement} es exponencialmente estable
para todo $\gamma>0$.
\end{proposicion}

\begin{proof}
La linealización es $\dot{\tilde p}=-\gamma E\tilde p$. Por el teorema de
Gershgorin, cada autovalor de $E$ está en un disco centrado en $E_{\ell\ell}>0$
de radio $\sum_{\ell'\neq\ell}|E_{\ell\ell'}|<E_{\ell\ell}$; ningún disco
contiene el origen ni corta el semiplano izquierdo, luego todos los autovalores
tienen parte real positiva y $-\gamma E$ es Hurwitz.
\end{proof}

\begin{observacion}[Cuándo se viola]
\label{obs:cascada}
La dominancia diagonal expresa que el efecto de un precio sobre su propia carga
domina su efecto sobre las demás. Se degrada cuando varias cargas comparten los
mismos AMR candidatos y sus urgencias suben a la vez: entonces las elasticidades
cruzadas crecen y el ajuste puede oscilar. El fenómeno no es hipotético; los
esquemas de tarificación por congestión desplegados en la práctica han exhibido
inestabilidades de este tipo. La consecuencia de diseño es que la condición debe
comprobarse sobre la matriz estimada, no suponerse, y que la saturación de la
urgencia es el mecanismo directo para preservarla.
\end{observacion}
